\subsection{在伯努利数上的应用}

VI, p.~286 的定理 1 表明，对 \(0 \leqslant x < \pi\) ，以 \(\frac{2x}{n^2\pi^2 - x^2} \geqslant 0\) 为通项的级数收敛。另一方面，对每个复数 \(z\) （满足 \(|z| < \pi\) 的），可以写成

\[
{\frac {2 z}{n ^ {2} \pi^ {2} - z ^ {2}}} = {\frac {2 z}{n ^ {2} \pi^ {2}}} \sum_ {k = 0} ^ {\infty} {\frac {z ^ {2 k}}{n ^ {2 k} \pi^ {2 k}}}
\]

右端的级数绝对收敛。我们将由此推出，“二重”级数

\[
\sum_ {n = 1} ^ {\infty} \sum_ {k = 1} ^ {\infty} \frac {- 2 z ^ {2 k - 1}}{n ^ {2 k} \pi^ {2 k}}\tag{11}
\]

在开圆盘 \(|z| < \pi\) 中绝对收敛，在含于此圆盘的每个紧集上正规收敛，且其和为 \(\cot z - \frac{1}{z}\) 。事实上，对 \(|z| \leqslant a < \pi\) ，(11) 的通项的绝对值至多等于 \(2a^{2k-1}/n^{2k}\pi^{2k}\) ，而项 \(2a^{2k-1}/n^{2k}\pi^{2k}\) 中任意有限多个之和小于有限数 \(\sum_{n=1}^{\infty} \frac{2a}{n^2\pi^2 - a^2}\) ；先对 \(k\) 求和、再对 \(n\) 求和，即可看出级数 (11) 的和等于 \(\sum_{n=1}^{\infty} \frac{2z}{z^2 - n^2\pi^2}\) ，这就证明了我们的结论。

现在若把级数 (11) 先对 \(n\) 求和、再对 \(k\) 求和，就得到恒等式（对 \(|z| < \pi\) ）

\[
\cot z - \frac {1}{z} = - 2 \sum_ {k = 1} ^ {\infty} \frac {S _ {2 k}}{\pi^ {2 k}} z ^ {2 k - 1}\tag{12}
\]

其中我们置 \(S_{k} = \sum_{n=1}^{\infty} \frac{1}{n^{k}}\) 。于是由 (1) （VI, p.~283）得到，对 \(|z| < 2\pi\)

\[
\frac {z}{e ^ {z} - 1} = 1 - \frac {z}{2} + \sum_ {n = 1} ^ {\infty} \frac {(- 1) ^ {n - 1} S _ {2 n}}{2 ^ {2 n - 1} \pi^ {2 n}} z ^ {2 n}\tag{13}
\]

由此得到公式

\[
b _ {2 n} = (- 1) ^ {n - 1} (2 n)! \frac {2 \mathrm{S} _ {2 n}}{(2 \pi) ^ {2 n}} \quad \text {   for   } n \geqslant 1,\tag{14}
\]

这个公式特别表明数 \(S_{2n} / \pi^{2n}\) 是有理数。显然 \(S_{k + 1} \leqslant S_k\) ，故对每个整数 \(k \geqslant 2\) 有 \(S_k \leqslant S_2 = \pi^2 / 6 \leqslant 2\) ；由 (14) 可推出伯努利数的下列不等式

\[
\frac {2 (2 n) !}{(2 \pi) ^ {2 n}} \leqslant | b _ {2 n} | \leqslant 4 \frac {(2 n) !}{(2 \pi) ^ {2 n}} \quad \text {   for   } n \geqslant 1.\tag{15}
\]

由这些不等式可以推出伯努利多项式 \(\mathrm{B}_n(x) = \sum_{k=0}^{n}\binom{n}{k}b_kx^{n - k}\) 的一个估计式；特别地，对 \(0 \leqslant x \leqslant 1\) 有

\[
\left| \mathrm{B} _ {n} (x) \right| \leqslant 4 \sum_ {k = 0} ^ {n} \binom {n} {k} \frac {k !}{(2 \pi) ^ {k}} = 4 \frac {n !}{(2 \pi) ^ {n}} \sum_ {k = 0} ^ {n} \frac {(2 \pi) ^ {k}}{k !} \leqslant 4 e ^ {2 \pi} \frac {n !}{(2 \pi) ^ {n}}.\tag{16}
\]

